\section*{第 \ref{ch:g11:ocean-ridges} 章——洋脊与大洋岩石圈}

\begin{solution}{exo:g11:ocean-ridges:1}
\omterm{def:g11:ocean-ridges:peridotite}{橄榄岩}是一种深色粗粒岩石，主要由橄榄石构成，并含少量
辉石；地幔的大部分由它组成。\omterm{def:g11:ocean-ridges:gabbro}{辉长岩}是一种深色粗粒火成岩，
由斜长石和辉石构成，是深部\omterm{def:g11:ocean-ridges:oceanic-lithosphere}{洋壳}中
最丰富的岩石。\omterm{def:g11:ocean-ridges:gabbro}{辉长岩}属于地壳，\omterm{def:g11:ocean-ridges:peridotite}{橄榄岩}属于地幔。
\omterm{def:g11:ocean-ridges:gabbro}{辉长岩}在海底深处缓慢结晶，岩浆冷却得很慢：
其晶体有足够时间长大。玄武岩熔岩在海底迅速冷却，
因而晶体仍然细小。
\end{solution}

\begin{solution}{exo:g11:ocean-ridges:2}
自海底向下依次为：沉积物（轴部无沉积物）、玄武岩熔岩（枕状
熔岩与席状流）、席状岩墙、\omterm{def:g11:ocean-ridges:gabbro}{辉长岩}，\omterm{def:g11:seismic-structure:moho}{莫霍面}以下则为
地幔的\omterm{def:g11:ocean-ridges:peridotite}{橄榄岩}。钻孔穿过 $1507/7100 \approx 0.21$，约占
地壳的五分之一：深部\omterm{def:g11:ocean-ridges:gabbro}{辉长岩}和\omterm{def:g11:seismic-structure:moho}{莫霍面}从未被钻探触及。
\end{solution}

\begin{solution}{exo:g11:ocean-ridges:3}
$d(0) = \qty{2500}{m}$; $d(9) = 2500 + 350 \times 3 = \qty{3550}{m}$;
$d(25) = 2500 + 350 \times 5 = \qty{4250}{m}$.
\end{solution}

\begin{solution}{exo:g11:ocean-ridges:4}
\omterm{def:g11:ocean-ridges:normal-fault}{正断层}是一种倾斜断层，沿其断层面，断层面以上的岩块
相对下方岩块向下移动。在慢速洋脊处，板块
将冷而脆的\omterm{def:g11:lithosphere-asthenosphere:lithosphere}{岩石圈}拉张：\omterm{def:g11:lithosphere-asthenosphere:lithosphere}{岩石圈}被拉长并断裂。
断层上方的岩块向洋脊轴部下滑，轴向谷地是
两组断层之间下降的岩块。
\end{solution}

\begin{solution}{exo:g11:ocean-ridges:5}
它不沸腾，原因在于压强：超过两千米的海水压在其上，
在此压强下，水在远高于
\qty{100}{\celsius}的温度下仍为液态。黑色的“烟”是由细小颗粒组成的云，
其中主要是硫化铁等矿物，它们在\omterm{def:g11:ocean-ridges:hydrothermal}{热液喷口}的热水
富含溶解金属，遇冷海水时沉淀而成。
\end{solution}

\begin{solution}{exo:g11:ocean-ridges:6}
（1）深约 \qty{54}{km} 处，温度约 \qty{1324}{\celsius}。（2）深约
\qty{85}{km} 处，温度约 $1400 + 0.45 \times 85 \approx
\qty{1438}{\celsius}$。（3）温度更高的地幔在更深处开始熔融，因此上升
的岩石柱中有更长一段在到达地壳前熔融，
生成更多玄武岩，地壳也更厚。
\end{solution}

\begin{solution}{exo:g11:ocean-ridges:7}
（1）$1489 - 1345 = \qty{144}{\celsius}$，低于其熔化温度。
（2）岩石上升 $100 - 54 = \qty{46}{km}$，二者之差
平均每千米缩小 $144/46 \approx \qty{3.1}{\celsius}$。（3）否：岩石上升时
甚至略有降温，约为 $0.45 \times 46 \approx
\qty{21}{\celsius}$。其熔化温度下降得快得多，
随压强减小而降低。
\end{solution}

\begin{solution}{exo:g11:ocean-ridges:8}
半扩张速率分别为每年 \qty{7.3}{cm} 和 \qty{1.65}{cm}，即
\qty{73}{km} 和 \qty{16.5}{km}（每百万年）。要在一侧形成 \qty{1000}{km} 的
洋底，东太平洋海隆需 $1000/73 \approx 13.7$ 百万年；
大西洋中脊需 $1000/16.5 \approx 61$ 百万年。
\end{solution}

\begin{solution}{exo:g11:ocean-ridges:9}
（1）\qty{5}{cm} 即 \qty{5e-5}{km}：$65000 \times 5 \times 10^{-5} =
\qty{3.25}{km^2}$ 的新洋底（每年）。（2）$3.25 \times 7.1 \approx
\qty{23}{km^3}$ 的地壳（每年）。（3）$3.6 \times 10^8 / 3.25 \approx 1.1
\times 10^8$ 年，约 1.1 亿年：按此速度，洋脊可以
在约一亿年内生成整个洋底，与
地球年龄相比很短。这就是没有哪段洋底很古老的原因。
\end{solution}

\begin{solution}{exo:g11:ocean-ridges:10}
（1）断层面与水平面夹角为 \ang{60}，下降量是
直角三角形的竖直边，其斜边为滑距：
$1.5/\sin\ang{60} \approx \qty{1.73}{km}$。（2）每条断层使两侧
水平地分开 $1.5/\tan\ang{60} \approx \qty{0.87}{km}$；两条
断层合计约 \qty{1.73}{km}。（3）否。\qty{25}{km} 是
下降岩块两侧断层之间的距离，而非拉伸
量，仅约 \qty{1.7}{km}；洋脊还
通过轴部新生的岩墙和熔岩不断张开，并非只靠断层。
\end{solution}

\begin{solution}{exo:g11:ocean-ridges:11}
$d(25) = 6400 - 3200\,e^{-25/62.8} \approx 6400 - 3200 \times 0.672 \approx
\qty{4251}{m}$; $d(100) = 6400 - 3200\,e^{-100/62.8} \approx 6400 - 3200
\times 0.203 \approx \qty{5749}{m}$。在 25 至 100 百万年之间，洋底
下沉约 \qty{1500}{m}。平方根定律给出
$2500 + 350 \times 5 = \qty{4250}{m}$ 即 25 百万年时的深度：两种
定律在二者均适用之处结果一致。
\end{solution}

\begin{solution}{exo:g11:ocean-ridges:12}
（1）$\alpha\,\Delta T\,L = 3.2 \times 10^{-5} \times 700 \times 100 =
\qty{2.24}{km}$。（2）$s = 3380 \times 2.24/(3380 - 1025) = 7571/2355 \approx
\qty{3.2}{km}$。（3）指数定律所允许的总下沉量为 $6400 - 2500 =
\qty{3900}{m}$。仅冷却和收缩便可解释其中约
五分之四，计算中的温降和
厚度均取简单数值。
\end{solution}

\begin{solution}{exo:g11:ocean-ridges:13}
（1）$54 - 7 = \qty{47}{km}$，对比 $85 - 7 = \qty{78}{km}$：较热的
柱体长 $78/47 \approx 1.66$ 倍。（2）$7.1 \times 1.66 \approx
\qty{11.8}{km}$。（3）实测的 \qty{10.3}{km} 比正常值厚，与
预测相符，但小于估算值：柱体最深部的数千米，
刚高于熔融起始深度，产生的熔体比浅部更少，因此
假定每千米熔体量相同，会高估地壳厚度。
\end{solution}

\begin{solution}{exo:g11:ocean-ridges:14}
（1）定律给出 \qty{3178}{m}、\qty{4160}{m}、\qty{4782}{m} 和
\qty{5267}{m}；实测中位数分别浅 \qty{64}{m}，
\qty{142}{m}、\qty{317}{m} 和 \qty{619}{m}。（2）差值随
年龄增大而增大。（3）定律给出的是洋底本身的深度，已对其沉积物作了修正；
实测的柱体带有沉积物，老洋底上的沉积物更厚，
并包括海山和火山高原，
它们都比周围的洋底更浅。
\end{solution}

\begin{solution}{exo:g11:ocean-ridges:15}
（1）$p = 1025 \times 9.8 \times 2500 \approx \qty{25.1}{MPa}$，加上
\qty{0.1}{MPa}，约为 \qty{25.2}{MPa}。（2）这高于临界
压强 \qty{22.064}{MPa}，喷口水无论多热都不会沸腾。
（3）压强达到 \qty{22.064}{MPa} 时的深度为
$(22.064 - 0.1) \times 10^6/(1025 \times 9.8) \approx \qty{2190}{m}$：凡
深于约 \qty{2200}{m} 的喷口都不会沸腾。
\end{solution}

\begin{solution}{pb:g11:ocean-ridges:1}
\textbf{1.} $14.6/2 = \qty{7.3}{cm}$ 每年，即每百万年
\qty{73}{km}。

\textbf{2.} 在洋脊轴部的深度处，$d(0) = \qty{2500}{m}$。

\textbf{3.} 自海底向下：玄武岩熔岩（枕状熔岩和席状流）、
席状岩墙与\omterm{def:g11:ocean-ridges:gabbro}{辉长岩}，地壳总厚约 \qty{7.1}{km}；在
\omterm{def:g11:seismic-structure:moho}{莫霍面}之下为地幔的\omterm{def:g11:ocean-ridges:peridotite}{橄榄岩}，一直延伸到\omterm{def:g11:lithosphere-asthenosphere:lithosphere}{岩石圈}底部，
而\omterm{def:g11:lithosphere-asthenosphere:lithosphere}{岩石圈}在洋脊轴部仍很薄。沉积物以后将在其上堆积。

\textbf{4.} 板块彼此远离，\omterm{def:g11:lithosphere-asthenosphere:lithosphere}{软流圈}在
洋脊下方上升。其温度每
千米仅降约 \qty{0.45}{\celsius}，而解链温度随压强
降低而急剧下降。约 \qty{54}{km} 深处两者相遇：岩石开始熔融，
并随上升继续熔融，无需外加热量（
\omterm{def:g11:ocean-ridges:partial-melting}{减压熔融}）。

\textbf{5.} 每百万年，洋脊两侧共新生 $2 \times 73 = \qty{146}{km}$
洋底；洋脊每千米长度的地壳体积为
$146 \times 7.1 \times 1 \approx \qty{1040}{km^3}$.

\textbf{6.} $73 \times t$：\qty{73}{km}、\qty{292}{km}、\qty{1168}{km} 和
\qty{4672}{km}.

\textbf{7.} $2000/73 \approx 27$，约 2700 万年。

\textbf{8.} 板块的速度与方向会在数千万
年的时间里改变，如夏威夷–帝王海山链的弯曲所示
（见\cref{ch:g11:plate-motions}）；今日的速率未必
能维持 6400 万年。

\textbf{9.} 远离轴部后，地壳不再位于岩浆透镜体之上：它
冷却了，其下方也不再有足够热量驱动
环流并加热海水。

\textbf{10.} 东太平洋海隆以西，该块体随
太平洋板块运动。太平洋板块向西以\omterm{def:g10:layered-ocean:thermocline}{深海}海沟为界，在那里
它俯冲于相邻板块之下：俯冲带（见\cref{ch:g11:subduction-zones}）。

\textbf{11.} $d(1) = \qty{2850}{m}$, $d(4) = 2500 + 700 = \qty{3200}{m}$,
$d(16) = 2500 + 1400 = \qty{3900}{m}$, $d(64) = 2500 + 2800 = \qty{5300}{m}$.

\textbf{12.} $3900 - 2500 = \qty{1400}{m}$，在最初的 1600 万年内，
而 $5300 - 3900 = \qty{1400}{m}$，对应 1600 万至 6400 万年之间：
下沉量相同，所用时间却长了三倍。洋底
随年龄增大，下沉越来越慢。

\textbf{13.} 由 $350\sqrt{t} = 1500$ 得 $\sqrt{t} \approx 4.29$，故
$t \approx 18.4$，约 1840 万年。

\textbf{14.} $6400 - 3200\,e^{-64/62.8} \approx 6400 - 3200 \times 0.361
\approx \qty{5245}{m}$，比平方根定律约低 \qty{55}{m}，两
者约在 \qty{1}{\percent} 以内一致。

\textbf{15.} 实测中位数比
定律约浅 \qty{650}{m}。定律描述的是已扣除沉积物重量的洋底；
实测深度包括沉积物，此类洋底上的沉积物已积累
数千万年，还包括海山和火山高原；它们都使
实测洋底变得更浅。

\textbf{16.} $\alpha\,\Delta T = 3.2 \times 10^{-5} \times 700 = 0.0224$：地幔
部分的密度增大 \qty{2.24}{\percent}，
$3380 \times 1.0224 \approx \qty{3456}{kg/m^3}$.

\textbf{17.} 地幔部分厚 $100 - 7.1 = \qty{92.9}{km}$。平均
密度：$(7.1 \times 2900 + 92.9 \times 3456)/100 \approx (20590 + 321060)/100
\approx \qty{3416}{kg/m^3}$.

\textbf{18.} $3416 > 3380$：整个老\omterm{def:g11:lithosphere-asthenosphere:lithosphere}{岩石圈}的密度
高出约 \qty{36}{kg/m^3}，约 \qty{1}{\percent}，比
其下方的\omterm{def:g11:lithosphere-asthenosphere:lithosphere}{软流圈}更致密。较轻的地壳已不足以
抵消其冷而致密的地幔部分。

\textbf{19.} 因其比\omterm{def:g11:lithosphere-asthenosphere:lithosphere}{软流圈}致密，它可以在
海沟处重新沉入地幔：俯冲，即下一章的主题。

\textbf{20.} 到 6400 万年时，该段洋底位于约
\qty{5300}{m}处，约 \textbf{5.3~千米}深，比其诞生时的位置深
\qty{2.8}{km}。
\end{solution}
